% Results-section addendum: interpretation and scope of the LCA estimates.
% Drafted 2026-05-13 for inclusion at the end of Section~\ref{sec:results}, after the hukou subsubsection.
% Standalone .tex; paste directly into sec_results.tex.
% Cross-references assume the labels in sec_model.tex and sec_results.tex resolve as written.

\subsection{What the LCA estimates identify}\label{sec:interpretation}

The estimates in this section rest on the joint identification assumptions A1--A5, of which A5 (the LCA restriction) is the one we test directly through Hansen's $J$-statistic.
Across our reported specifications, the test fails to reject at conventional levels in Indonesia and Tanzania, as well as in both hukou-split subsamples for China.
A2 (the i.i.d.\ assumption on $\nu_{it}^l$) is maintained rather than tested.
Its role inside the identification set is to underwrite the reading that trajectory membership reflects comparative advantage cleanly.
The estimated $\phi$, $\beta$, $\mu_{\underline{d}}$, and trajectory-specific $\Delta_{\underline{d}}$ are conditional averages of $\Delta_i$ across trajectories, identified by the GMM moment conditions in equation~\eqref{eq:gmm-moments} under A1--A5 jointly.

Read against the model's decision rule, the estimates nonetheless suggest that observed sorting does not reflect comparative advantage alone.
Under negative $\phi$ and the rule in equation~\eqref{eq:decision-rule}, the always-urban trajectory $d_T$ should aggregate workers with high pecuniary returns to migration.
The restricted GRC instead places $\Delta_{d_T}$ near zero or below zero in Indonesia and Tanzania.
A symmetric pattern holds for never-migrants: under rational sorting the model would predict $\Delta_{d_N}$ near zero or negative, while our estimates place it consistently well above zero.

We do not read this pattern as evidence against A5.
The LCA restriction holds in the $J$-test sense, and the conditional averages it identifies are well-defined population objects.
What the pattern instead suggests is that A2 (the i.i.d.\ assumption on $\nu_{it}^l$) is strained: trajectories evidently reflect features beyond comparative advantage alone.
Persistent amenity value, social ties, and costs of return migration can keep workers in urban locations whose pecuniary returns to migration are not large.
Symmetric frictions can keep workers rural whose pecuniary returns to migration would be substantial if realized.
Section~\ref{sec:model_identification} already concedes this possibility; our estimates put magnitudes on it.
The testable consequence of A2 being strained would be a departure from A5, which the $J$-test would detect.
That this does not happen in our reported specifications means the strain on A2 does not break the LCA conditional-mean structure.
The estimated $\Delta_{\underline{d}}$ remain consistent estimates of trajectory-conditional averages; what loosens is the structural reading that those averages represent what rational sorters would have realized.

The China hukou comparison sharpens the reading.
Where the institutional barrier is absent, in the urban hukou subsample, the LCA estimates align with rational sorting under the decision rule: the estimated return for never-migrants is close to zero, consistent with workers who chose rural because their personal pecuniary return to urban location is approximately zero.
Where the barrier binds, in the rural hukou subsample, the alignment breaks: the estimated return for never-migrants is large and positive, consistent with workers whose pecuniary returns to migration are unrealized.
The contrast in $\phi$ across regimes is direct evidence that institutional context shapes which workers reach which trajectories, separately from comparative advantage.

The distinction between identification and behavioral interpretation also clarifies how to read counterfactual statements built from our estimates.
The LCA structure identifies the consumption side of $\Delta_i$.
Reallocations that would place workers in locations they currently avoid would raise consumption to the extent characterized by our estimates, but would also forgo non-pecuniary value not captured in $\Delta_i$.
Our estimates are therefore well-suited to consumption-side comparisons; they are conservative when extended to welfare more broadly.
